% -----------------------------------------------------------
\section{Christ (Jesus)}
% -----------------------------------------------------------
	\label{CHRIST}
	
	This entry is for discussion of the individual called Jesus Christ.
	
% % ----------------------
% \subsection{See also}
% % ----------------------

% % ----------------------
% \subsection{1828 American Dictionary of the English Language} % *1828
% % ----------------------
	% \label{CHRIST-1828}

	% \begin{compactitem}
		% \item
	% \end{compactitem}

% % ----------------------
% \subsection{Oxford English Dictionary} % *OED
% % ----------------------
	% \label{CHRIST-OED}

	% \begin{compactitem}
		% \item[]
	% \end{compactitem}

% ----------------------
\subsection{Scriptural Usage} % *SCRIPTURES
% ----------------------
	\label{CHRIST-SCRIPTURES}

	% % -----
	% \subsubsection{Old Testament} % *OT
	% % -----
		% \label{CHRIST-OT}
	
		% % --
		% \paragraph{KJV}
		% % --
			% \label{CHRIST-OT-KJV}
	
			% \begin{tabular}{ll}
				% Total occurrences in the OT & 
			% \end{tabular}
			
			% \begin{compactdesc}
				% \item[]
			% \end{compactdesc}
			
		% % --
		% \paragraph{Hebrew Bible}
		% % --
			% \label{CHRIST-HEBREW}
		
			% %
			% \subparagraph{\texorpdfstring{\he{sample word}}{}} % paste actual hebrew text here
			% %
				% \label{PAGA} % whatever is in step bible without periods
			
				% \begin{compactdesc}
					% \item[HALOT , PDF ] \textbf{} % Use the 1994 PDF
						% \begin{compactitem}
							% \item[1]
						% \end{compactitem}
					% \item[BDB , PDF ] \textbf{}
						% \begin{compactitem}
							% \item[1]
						% \end{compactitem}
					% \item[DCH PDF ] \textbf{} % don't include the actual page number, they repeat from volume to volume
						% \begin{compactitem}
							% \item[1]
						% \end{compactitem}
				% \end{compactdesc}
				
				% \begin{compactdesc}
					% \item[Some OT uses] \textbf{}
						% \begin{compactdesc}
							% \item[]
						% \end{compactdesc}
				% \end{compactdesc}
			
		% % --
		% \paragraph{Septuagint}
		% % --
			% \label{CHRIST-LXX}
		
			% %
			% \subparagraph{\texorpdfstring{\gr{sample word}}{}}
			% %
		
	% -----
	\subsubsection{New Testament} % *NT
	% -----
		\label{CHRIST-NT}
	
		% % --
		% \paragraph{KJV}
		% % --
			% \label{CHRIST-NT-KJV}		
	
			% \begin{tabular}{ll}
				% Total occurrences in the NT & 
			% \end{tabular}
			
			% \begin{compactdesc}
				% \item[]
			% \end{compactdesc}
			
		% --
		\paragraph{Greek New Testament} % *GREEK
		% --
			\label{CHRIST-GREEK}
		
			%
			\subparagraph{\texorpdfstring{\gr{Χριστός}}{}}
			%
				\label{CHRISTOS}
			
				\begin{compactdesc}
					\item[BDAG 971b] \textbf{}
						\begin{compactitem}
							\item[1] \textbf{fulfiller of Israelite expectation of a deliverer, \textit{the Anointed One, the Messiah, the Christ}}, appellative...\gr{ἐπυνθάνετο ποῦ ὁ Χριστὸς γεννᾶται} \textit{he inquired where the Messiah was to be born} (\hyperref[MATTHEW-2-4]{Matthew 2:4})...etc. The transition to sense 2 [below] is marked by certain passages in which \gr{Χριστός} does not mean the Messiah in general (even when the reference is to Jesus), but a very definite Messiah, Jesus, who is now called \textit{Christ} not as a title but as a name
							\item[2] \textbf{the personal name ascribed to Jesus, \textit{Christ}}, which many gentiles must have understood in this way (to them it seemed very much like \gr{Χρηστός}...a name that is found in literature
						\end{compactitem}
					% \item[CGL , PDF ] \textbf{}
						% \begin{compactitem}
							% \item 
						% \end{compactitem}
					% \item[LSJ , PDF ] \textbf{}
						% \begin{compactitem}
							% \item 
						% \end{compactitem}
				\end{compactdesc}
				
				% \begin{tabular}{ll}
					% Total occurrences in the NT & 
				% \end{tabular}
				
				% \begin{compactdesc}
					% \item[Some NT uses] \textbf{}
						% \begin{compactdesc}
							% \item[]
						% \end{compactdesc}
				% \end{compactdesc}
		
	% % -----
	% \subsubsection{Book of Mormon} % *BOM
	% % -----
		% \label{CHRIST-BOM}
	
		% \begin{tabular}{ll}
			% Total occurrences in the BOM & 
		% \end{tabular}
		
		% \begin{compactdesc}
			% \item[]
		% \end{compactdesc}
		
	% % -----
	% \subsubsection{Doctrine and Covenants} % *DC
	% % -----
		% \label{CHRIST-DC}
	
		% \begin{tabular}{ll}
			% Total occurrences in the DC & 
		% \end{tabular}
		
		% \begin{compactdesc}
			% \item[]
		% \end{compactdesc}
		
	% % -----
	% \subsubsection{Pearl of Great Price} % *PGP
	% % -----
		% \label{CHRIST-PGP}
	
		% \begin{tabular}{ll}
			% Total occurrences in the PGP & 
		% \end{tabular}
		
		% \begin{compactdesc}
			% \item[]
		% \end{compactdesc}
		
% ----------------------
\subsection{Key Phrases} % *KEYPHRASES *PHRASES
% ----------------------
	\label{CHRIST-PHRASES}

	\begin{compactdesc}
	
		\item[light of Christ] See \hyperref[LIGHT-PHRASES]{here}.
		
		\item[spirit of Christ] See \hyperref[SPIRIT-PHRASES]{here}.
		
		% \item[] \textbf{\#times} | 
			% \begin{compactdesc}
				% \item[Bible anchor verses] \phantom{.}
					% \begin{compactdesc}
						% \item[Romans 5:5] \gr{}
					% \end{compactdesc}
				% \item[Commentary] ---
			% \end{compactdesc}
			
	\end{compactdesc}
	
% % ----------------------
% \subsection{Bible Dictionaries} % *DICTIONARIES
% % ----------------------
	% \label{CHRIST-BD}

% % ----------------------
% \subsection{Restoration Usage} % *RESTORATION
% % ----------------------
	% \label{CHRIST-RESTORATION}

	% % -----
	% \subsubsection{Joseph Smith} % *JS
	% % -----
		% \label{CHRIST-JS}
	
		% \begin{compactdesc}
			% \item[{\hyperref[KEY]{Date/Source}}]
		% \end{compactdesc}
		
	% % -----
	% \subsubsection{Temple Ordinances}
	% % -----
		% \label{CHRIST-TEMPLE}
	
		% \begin{compactdesc}
			% \item[{\hyperref[KEY]{Date/Source}}]
		% \end{compactdesc}
	
	% % -----
	% \subsubsection{Brigham Young} % *BY
	% % -----
		% \label{CHRIST-BY}
	
		% \begin{compactdesc}
			% \item[{\hyperref[KEY]{Date/Source}}]
		% \end{compactdesc}
	
% ----------------------
\subsection{Commentary and Studies} % *COMMENTARY *STUDIES
% ----------------------
	\label{CHRIST-COMMENTARY}

	% % -----
	% \subsubsection{Key Points}
	% % -----
		% \label{CHRIST-KEYS}
	
		% \begin{compactitem}
			% \item
		% \end{compactitem}
		
	% -----
	\subsubsection{Christ progressing}
	% -----
		\label{CHRIST-progressing}
		
		\begin{compactitem}
			\item \hyperref[PHILIPPIANS-2-5]{Philippians 2:5--11}
				\begin{compactitem}
					\item Potentially, I'm not sure.
				\end{compactitem}
			\item \hyperref[DC93]{DC 93}, many parts of it, especially from the beginning up to and including verse 22.
			\item \hyperref[3NEPHI-12-48]{3 Nephi 12:48}, with the difference to Matthew 5:48, in that Christ is now also perfect.
			\item Orson Pratt, 28 August 1852, Deseret News Extra, 16a, first new paragraph
		\end{compactitem}
		
	% -----
	\subsubsection{Titles for Jesus Christ}
	% -----
		\label{CHRIST-titles}
	
		\begin{compactdesc}
			\item[advocate] See \hyperref[DC45-3]{DC 45:3} and my notes there.
			\item[Alpha/Omega] See \hyperref[REVELATIONNT-1-8]{Revelation 1:8} and \hyperref[DC95-17]{DC 95:17}. Could be referring to eternal lives/\hyperref[ROUND-PHRASES]{eternal round} stuff.
			\item[mediator] See \hyperref[HEBREWS-12-24]{Hebrews 12:24} and \hyperref[DC76-69]{DC 76:69}.
			\item[mighty one of Israel] See \hyperref[ISAIAH-1-24]{Isaiah 1:24}, \hyperref[ISAIAH-30-29]{Isaiah 30:29}, \hyperref[1NEPHI-22-12]{1 Nephi 22:12}, and \hyperref[DC36-1]{DC 36:1}
			\item[only begotten] I understand this term to refer to Christ as the only individual who receives life and power directly through Heavenly Father; we receive it through Christ. See \hyperref[MOSES-4-1]{Moses 4:1} and notes around there.
				\begin{compactitem}
					\item Similarly, see \hyperref[DC93-12]{DC 93:12--14}, along with \hyperref[MOSIAH-15-1]{Mosiah 15:1--5}.
				\end{compactitem}
			\item[Lamb of God] \hyperref[1NEPHI-11-21]{1 Nephi 11:21}
			\item[Eternal Father] \hyperref[1NEPHI-11-21]{1 Nephi 11:21}
			\item[Everlasting God] \hyperref[1NEPHI-11-32]{1 Nephi 11:32}
		\end{compactdesc}
		
	% -----
	\subsubsection{The in-dwelling relationship of Christ and the Father, and Christ and us}
	% -----
		\label{CHRIST-indwelling}
		
		See:
		
			\begin{compactitem}
				\item \hyperref[JOHN-15-1]{John 15:1--11}
				\item \hyperref[JOHN-17]{John 17}
					\begin{compactitem}
						\item Not the entire chapter, but it does appear several times.
					\end{compactitem}
				\item \hyperref[2CORINTIHANS-5]{2 Corinthians 5}
					\begin{compactitem}
						\item The idea occurs a few times in this chapter.
					\end{compactitem}
				\item \hyperref[GALATIANS-2-20]{Galatians 2:20}
				\item \hyperref[GALATIANS-3-26]{Galatians 3:26--29}
				\item \hyperref[EPHESIANS-1-10]{Ephesians 1:10--11}
				\item \hyperref[PHILIPPIANS-3-8]{Philippians 3:8--14}
				\item \hyperref[MOSIAH-15-1]{Mosiah 15:1--5}
				\item \hyperref[3NEPHI-1-14]{3 Nephi 1:14}
				\item \hyperref[3NEPHI-9-15]{3 Nephi 9:15}
				\item \hyperref[3NEPHI-11-27]{3 Nephi 11:27}
				\item \hyperref[3NEPHI-19-23]{3 Nephi 19:23}
				\item \hyperref[3NEPHI-19-29]{3 Nephi 19:29}
				\item \hyperref[3NEPHI-28-10]{3 Nephi 28:10}
				\item \hyperref[DC35-2]{DC 35:2}
					\begin{compactitem}
						\item See my notes here, I have a lot of connections to other verses.
					\end{compactitem}
				\item \hyperref[DC93-16]{DC 93:16--17}
			\end{compactitem}
		
	% -----
	\subsubsection{Christ filling a certain role in the plan of salvation}
	% -----
		\label{CHRIST-role}
		
		\begin{compactitem}
			\item \hyperref[2CORINTHIANS-1-19]{2 Corinthians 1:19--20}
			\item \hyperref[DC13]{DC 13} and my note on ``Messiah''
			\item \hyperref[DC18-21]{DC 18:21--25}, which is a discussion of the name of Christ, but especially verses 23--24 talk about the name in an unusual way I think
		\end{compactitem}
		
	% -----
	\subsubsection{Christ experiencing temptation}
	% -----
	
		See \hyperref[TEMPTATION-christ]{this study}.
		
	% -----
	\subsubsection{Other aspects of Christ's character}
	% -----
		\label{CHRIST-character}
		
		\begin{compactitem}
			\item \hyperref[PHILIPPIANS-2-5]{Philippians 2:5--11}
			\item \hyperref[COLOSSIANS-1-9]{Colossians 1:9--29}
				\begin{compactitem}
					\item These verses go with places like DC 93 and DC 88, in giving us very fundamental knowledge about Christ's character.
				\end{compactitem}
			\item \hyperref[1NEPHI-11-18]{1 Nephi 11:18}, Mary is the mother of Christ ``after the manner of the flesh''---but not in some other way
			\item \hyperref[JAROM-1-11]{Jarom 1:11}---something really interesting about the temporal nature of Christ and the atonement. There is another verse, maybe in Jacob or 2 Nephi, that says similar things, I think.
			\item \hyperref[MOSIAH-7-27]{Mosiah 7:27}
			\item \hyperref[MOSIAH-13-34]{Mosiah 13:34}
			\item \hyperref[MOSIAH-15-1]{Mosiah 15:1--5}
		\end{compactitem}